\section{Extended Experimental Contract}
\label{app:experiment_contract}

This appendix records analyses that are required for the final evidence bundle
but may move between main paper and supplementary material after results and
the official CVPR 2027 format are frozen.

\subsection{Artifact Manifest}
\ClaimMap{C10,C11,C12,C15}

Every table-eligible run must provide: model commit; evaluator commit;
configuration and preprocessing fingerprints; seed; hardware and software
environment; MultiRep release, source, download time, split, and checksums;
per-video and per-identity ground truth and predictions; oracle/predicted-track
status; logs; and generated-table provenance. An artifact lacking any required
field is retained for diagnosis but cannot support a paper claim. UCFRep's
sealed test partition remains untouched until the method, checkpoints, and
analysis code are frozen.

\begin{table*}[t]
  \centering
  \scriptsize
  \setlength{\tabcolsep}{1.7pt}
  \caption{Required characterization of the frozen MultiRep split before
  hypothesis testing. Every tempo or corruption stratum must also report its
  exact number of videos, people, and events.}
  \label{tab:dataset_characterization}
  \begin{tabular}{lcccccccc}
    \toprule
    Split & Videos & People & Events & Active IDs & Tempo-change &
    Pauses & Gaps & Window-eligible \\
    \midrule
    Train & \Protocol{data.train.videos} & \Protocol{data.train.people} &
      \Protocol{data.train.events} & \Protocol{data.train.active_ids} &
      \Protocol{data.train.tempo_tracks} & \Protocol{data.train.pauses} &
      \Protocol{data.train.gaps} & \Protocol{data.train.window_eligible} \\
    Validation & \Protocol{data.val.videos} & \Protocol{data.val.people} &
      \Protocol{data.val.events} & \Protocol{data.val.active_ids} &
      \Protocol{data.val.tempo_tracks} & \Protocol{data.val.pauses} &
      \Protocol{data.val.gaps} & \Protocol{data.val.window_eligible} \\
    Test & \Protocol{data.test.videos} & \Protocol{data.test.people} &
      \Protocol{data.test.events} & \Protocol{data.test.active_ids} &
      \Protocol{data.test.tempo_tracks} & \Protocol{data.test.pauses} &
      \Protocol{data.test.gaps} & \Protocol{data.test.window_eligible} \\
    \bottomrule
  \end{tabular}
\end{table*}

\begin{table*}[t]
  \centering
  \footnotesize
  \setlength{\tabcolsep}{2.5pt}
  \caption{Scope matrix for the three closest method families and this
  pre-results specification. ``Intended'' and ``provisional'' entries are not
  implementation or performance claims.}
  \label{tab:prior_scope}
  \begin{tabular}{p{0.13\textwidth}p{0.12\textwidth}p{0.14\textwidth}
    p{0.18\textwidth}p{0.14\textwidth}p{0.18\textwidth}}
    \toprule
    Method & Counting input & Output scope & Counting supervision & Tempo unit & Identity mechanism \\
    \midrule
    PAMS~\cite{gao2026pams} & Pose stream & Single stream & Self-supervised representation & Period-adaptive scales & None \\
    MultiCounter~\cite{tang2024multicounter} & RGB video & Person/events & Supervised MRAC & Per-person intervals & Joint association \\
    MultiCounter+~\cite{tang2026multicounterplus} & RGB video & Person/events & Supervised MRAC & Long/short-aware & Dynamic Siamese query \\
    Ours (specification) & Supplied pose tracks & Oracle person / predicted track & Intended no person count/period & Window-local routing & Supplied track index \\
    \bottomrule
  \end{tabular}
\end{table*}

\subsection{Mechanism Ablations}
\ClaimMap{C03,C04,C05,C06,C12}

The mechanism study compares one global tempo per track against window-local
tempo; uniform, hard, and soft routing; one, two, three, and five experts;
PAMS independently applied to each track; and sliding-window PAMS with
deterministic tempo bins. Boundary alternatives include non-overlapping
windows, center-crop decoding, uniform response averaging, phase-aligned
fusion, normalized overlap-add, and integer-window summation as a diagnostic
failure mode. Training is compared with and without the synchronized
within-track continuity term. An oracle router based on ground-truth periods
may serve as an upper-bound diagnostic only when the required annotation
exists, and must never be grouped with deployable methods. Arbitrary
person-indexed parameter sets are excluded because identities do not persist
across videos; efficiency instead compares shared batched execution with
independent replicas of the same frozen checkpoint.

\begin{table*}[t]
  \centering
  \footnotesize
  \setlength{\tabcolsep}{3pt}
  \caption{Planned mechanism analysis. Parameter counts expose the cost of
  expert count and weight sharing.}
  \label{tab:mechanisms}
  \begin{tabular}{llcccc}
    \toprule
    Factor & Setting & Period-mAP$\uparrow$ & AvgMAE$\downarrow$ &
    AvgOBO$\uparrow$ & Params (M)$\downarrow$ \\
    \midrule
    Tempo scope & Track-global & \R{mech.tempo.global.pmap} &
      \R{mech.tempo.global.mae} & \R{mech.tempo.global.obo} &
      \R{mech.tempo.global.params} \\
    & Window-local & \R{mech.tempo.local.pmap} &
      \R{mech.tempo.local.mae} & \R{mech.tempo.local.obo} &
      \R{mech.tempo.local.params} \\
    Router & Uniform & \R{mech.router.uniform.pmap} &
      \R{mech.router.uniform.mae} & \R{mech.router.uniform.obo} &
      \R{mech.router.uniform.params} \\
    & Hard & \R{mech.router.hard.pmap} & \R{mech.router.hard.mae} &
      \R{mech.router.hard.obo} & \R{mech.router.hard.params} \\
    & Soft & \R{mech.router.soft.pmap} & \R{mech.router.soft.mae} &
      \R{mech.router.soft.obo} & \R{mech.router.soft.params} \\
    Boundary & Center crop & \R{mech.boundary.crop.pmap} &
      \R{mech.boundary.crop.mae} & \R{mech.boundary.crop.obo} &
      \R{mech.boundary.crop.params} \\
    & Response average & \R{mech.boundary.average.pmap} &
      \R{mech.boundary.average.mae} & \R{mech.boundary.average.obo} &
      \R{mech.boundary.average.params} \\
    & Normalized overlap-add & \R{mech.boundary.ola.pmap} &
      \R{mech.boundary.ola.mae} & \R{mech.boundary.ola.obo} &
      \R{mech.boundary.ola.params} \\
    \bottomrule
  \end{tabular}
\end{table*}

\subsection{Robustness and Tracking Diagnostics}
\ClaimMap{C02,C05,C12,C15}

Natural tempo variation will be stratified by within-track cycle-duration
dispersion. Deterministic count-preserving piecewise time warps, contiguous
masking, and injected identity switches will measure degradation from each
clean input using identical corruptions across methods. Perturbations are
separated by pipeline level. Pre-tracking video corruptions rerun the frozen
frontend and may change HOTA, IDF1, and IDSW; post-track pose/mask corruptions
hold tracks fixed, so tracking metrics must remain unchanged. The protocol
must record whether an identity switch changes only the track identifier or
also the pose content. Results will be stratified by active-person count,
asynchronous starts and stops, pauses, occlusion, and track fragmentation.

\begin{table*}[t]
  \centering
  \small
  \caption{Planned robustness report. Parenthesized values will be paired
  changes from the corresponding clean input with video-cluster bootstrap
  intervals.}
  \label{tab:robustness}
  \begin{tabular}{llllcccc}
    \toprule
    Condition & Level & Severity & Stratum \(n\) & Period-mAP$\uparrow$ & AvgMAE$\downarrow$ &
    AvgOBO$\uparrow$ & IDF1$\uparrow$ \\
    \midrule
    Clean & -- & -- & \Protocol{robust.clean_n} & \R{robust.clean.pmap} & \R{robust.clean.mae} &
      \R{robust.clean.obo} & \R{robust.clean.idf1} \\
    Natural tempo change & Data stratum & Low/medium/high & \Protocol{robust.tempo_n} & \R{robust.tempo.pmap} &
      \R{robust.tempo.mae} & \R{robust.tempo.obo} & \R{robust.tempo.idf1} \\
    Piecewise time warp & Post-track & \Protocol{robust.warp_levels} & \Protocol{robust.warp_n} & \R{robust.warp.pmap} &
      \R{robust.warp.mae} & \R{robust.warp.obo} & \Protocol{robust.fixed_idf1} \\
    Contiguous occlusion & Pre-track & \Protocol{robust.occlusion_levels} &
      \Protocol{robust.occlusion_n} & \R{robust.occlusion.pmap} & \R{robust.occlusion.mae} &
      \R{robust.occlusion.obo} & \R{robust.occlusion.idf1} \\
    Injected ID switches & Post-track & \Protocol{robust.switch_levels} &
      \Protocol{robust.switch_n} & \R{robust.switch.pmap} & \R{robust.switch.mae} &
      \R{robust.switch.obo} & \R{robust.switch.idf1} \\
    \bottomrule
  \end{tabular}
\end{table*}

\subsection{Generalization and Efficiency}
\ClaimMap{C12,C15}

If protocol-compatible, a model trained on MultiRep will be evaluated without
target-domain adaptation on the official RepCount and UCFRep single-person
protocols. Dataset-native normalized MAE and OBO will retain their original
names rather than being relabeled as MultiRep AvgMAE/AvgOBO. Such transfer can
support representation generalization only, not MRAC identity robustness.
Efficiency will
separate detector/pose/tracker preprocessing from counting latency and report
parameters, FLOPs, peak memory, latency, counting-only throughput, and
full-pipeline throughput as active identities increase. Shared batched
execution and independent same-checkpoint replicas are timed separately. If
ByteTrack~\cite{zhang2022bytetrack} is selected as the
common tracker, its exact code and checkpoint revision will be frozen in the
manifest; this mention does not pre-select the frontend.

\begin{table*}[t]
  \centering
  \small
  \caption{Planned efficiency report under one frozen hardware and timing
  protocol. The final generator expands every method into one row for each
  integer active-ID level from one through the frozen dataset maximum; the
  current Active IDs cell marks that pending grid rather than one aggregate
  measurement.}
  \label{tab:efficiency}
  \begin{tabular}{@{}lccccccc@{}}
    \toprule
    Method & Active IDs & Params (M)$\downarrow$ & GFLOPs$\downarrow$ & Peak GB$\downarrow$ &
    Preprocess ms/frame$\downarrow$ & Count ms/frame$\downarrow$ & Pipeline FPS$\uparrow$ \\
    \midrule
    MultiCounter & \Protocol{eff.identity_levels} & \R{eff.multicounter.params} & \R{eff.multicounter.flops} &
      \R{eff.multicounter.vram} & \R{eff.multicounter.pre} &
      \R{eff.multicounter.model} & \R{eff.multicounter.fps} \\
    MultiCounter+ & \Protocol{eff.identity_levels} & \R{eff.multicounterplus.params} &
      \R{eff.multicounterplus.flops} & \R{eff.multicounterplus.vram} &
      \R{eff.multicounterplus.pre} & \R{eff.multicounterplus.model} &
      \R{eff.multicounterplus.fps} \\
    \textbf{Ours} & \Protocol{eff.identity_levels} & \R{eff.ours.params} & \R{eff.ours.flops} &
      \R{eff.ours.vram} & \R{eff.ours.pre} & \R{eff.ours.model} &
      \R{eff.ours.fps} \\
    \bottomrule
  \end{tabular}
\end{table*}

Qualitative figures must be generated directly from frozen per-frame artifacts
and include persistent identity colors, valid masks, local period estimates,
routing weights, fused responses, and ground-truth/predicted boundaries. At
least one failure case is mandatory; curves may not be manually redrawn.
